%The mostly commonly used tool for forced alignment is the Montreal Forced Aligner, which wraps Hybrid HMM/DNN implemented in \textsc{kaldi} into an easy to use python interface. It does not  handle speaker attribution, and generally fails on overlapped speech as will be apparent in Section \ref{sec:experiments}.

%While there has been very little work on multi-talker Viterbi forced alignment apart from our prior work on shuffle models \cite{wiesner2026modeling}, a significant body of work on long-form alignment exists. 

%Recent work \cite{zeyer2026gradient} shows how any ASR system through which we can backpropagate can be repurposed for forced alignment by doing dynamic time warping (DTW) on gradients of the input signal to the system.

%We will have to talk about some systems that must exist for alignment from separation systems. 
%We probably also need to talk about some approaches to long-form alignment.

Forced alignment of single-speaker audio is a well-established problem. HMM-based aligners such as MFA~\cite{mcauliffe2017montreal, povey2011kaldi}, CTC-based aligners~\cite{CTC_align, rastorgueva2023nfa, bain2023whisperx} and recent novel approaches based on saliency~\cite{zeyer2026gradient}, all model one token sequence and one active speaker per frame, and for long recordings they avoid the full Viterbi lattice by segmenting the audio with voice activity detection (VAD) or a first decoding pass and aligning each segment separately~\cite{moreno1998recursive, katsamanis2011sailalign, biased_LMs}. This strategy requires either silent regions at which to segment or a first pass that is accurate enough to provide anchors. The alternative is to keep the exact search and reduce its memory, for instance by checkpointing the Viterbi lattice~\cite{tarnas1998reduced} or by pruning it to a band~\cite{kuang2022pruned}; the techniques of Section~\ref{sec:longform} belong to this second approach.

On the multi-talker side, the situation is reversed: several systems require aligned data but no current single technique can produce it. For example, token-level serialized output training~\cite{kanda2020sot, kanda22_interspeech} is trained on timed, speaker-attributed transcripts, which are currently obtained by aligning the single-speaker source recordings before they are mixed~\cite{polok2026mindgap}. Large-scale data pipelines such as Emilia~\cite{he2024emilia} and Sommelier~\cite{jung2026sommelier} instead apply source separation and diarization to obtain single-speaker segments that a standard aligner can process. This is effective on web audio with limited overlap, but in far-field conversational settings neural speech separation degrades and ceases to be effective, in particular as a front-end for ASR~\cite{cornell2025recent}. Apart from our previous work~\cite{wiesner2026modeling}, we are not aware of prior work on Viterbi forced alignment of overlapped multi-talker transcripts. Moreover, word timing alone is not sufficient for most applications: to be useful, multi-talker alignment must also attribute each word to a speaker. Attributing utterances to speakers without labels is an instance of the speaker partitioning problem of Br\"ummer and de Villiers~\cite{brummer2010partitioning}, whose solution space grows with the Bell numbers and which clustering methods such as VBx~\cite{landini2022vbx} solve approximately. The method of Section~\ref{sec:attribution} differs from both in that it scores speaker partitions of the reference utterances directly under the acoustic model, given a fixed token alignment.